\documentclass[twocolumn]{aastex631}
% Auto-generated by jansky_research.dmzcal.write_paper -- do not edit.
% dmzReal* from results/dmzcal_{provenance,metrics}.json (real data);
% dmzSyn* from results/dmzcal_{controls,dmz_power}.json (synthetic controls).
\newcommand{\dmzRealSource}{real: plan 99 coverage audit of E1-E3 + E4 point accuracy on FRBs/FRB localized hosts (certification sample per dmzcal_provenance.json)}
\newcommand{\dmzSynSource}{synthetic controls: plan 99 C0 (null false-fail) and C1 (power), E2 generative model only -- no real burst used}
\newcommand{\dmzRealNLocal}{105}
\newcommand{\dmzRealNProd}{48}
\newcommand{\dmzRealNCert}{36}
\newcommand{\dmzRealNChime}{20}
\newcommand{\dmzRealNDsa}{11}
\newcommand{\dmzRealNAskap}{4}
\newcommand{\dmzRealNMeer}{1}
\newcommand{\dmzRealEOneCovSixtyEight}{0.56}
\newcommand{\dmzRealEOneCovNinetyFive}{0.81}
\newcommand{\dmzRealEOneMedPit}{0.185}
\newcommand{\dmzRealEOneKsP}{1\times10^{-4}}
\newcommand{\dmzRealEOneTailLo}{0.42}
\newcommand{\dmzRealEOneTailHi}{0.03}
\newcommand{\dmzRealFrbOneMed}{2.75}
\newcommand{\dmzRealFrbOneLoSixtyEight}{2.04}
\newcommand{\dmzRealFrbOneHiSixtyEight}{3.18}
\newcommand{\dmzRealETwoCovSixtyEight}{0.22}
\newcommand{\dmzRealETwoCovNinetyFive}{0.50}
\newcommand{\dmzRealETwoMedPit}{0.027}
\newcommand{\dmzRealETwoKsP}{6\times10^{-13}}
\newcommand{\dmzRealETwoTailLo}{0.75}
\newcommand{\dmzRealETwoTailHi}{0.03}
\newcommand{\dmzRealFrbTwoMed}{3.16}
\newcommand{\dmzRealFrbTwoLoSixtyEight}{2.60}
\newcommand{\dmzRealFrbTwoHiSixtyEight}{3.57}
\newcommand{\dmzRealEThreeCovSixtyEight}{0.17}
\newcommand{\dmzRealEThreeCovNinetyFive}{0.44}
\newcommand{\dmzRealEThreeMedPit}{0.016}
\newcommand{\dmzRealEThreeKsP}{9\times10^{-20}}
\newcommand{\dmzRealEThreeTailLo}{0.83}
\newcommand{\dmzRealEThreeTailHi}{0.00}
\newcommand{\dmzRealFrbThreeMed}{2.42}
\newcommand{\dmzRealFrbThreeLoSixtyEight}{2.16}
\newcommand{\dmzRealFrbThreeHiSixtyEight}{2.64}
\newcommand{\dmzRealEOneNBelowHalf}{28}
\newcommand{\dmzRealDzOne}{+0.06}
\newcommand{\dmzRealDzTwo}{+0.12}
\newcommand{\dmzRealDzThree}{+0.15}
\newcommand{\dmzRealDzFour}{+0.06}
\newcommand{\dmzRealDzFourN}{35}
\newcommand{\dmzRealDmzOneKsP}{0.78}
\newcommand{\dmzRealDmzOneCovSixtyEight}{0.69}
\newcommand{\dmzRealDmzOneCovNinetyFive}{0.94}
\newcommand{\dmzRealDmzTwoKsP}{0.63}
\newcommand{\dmzRealDmzTwoCovSixtyEight}{0.75}
\newcommand{\dmzRealDmzTwoCovNinetyFive}{0.94}
\newcommand{\dmzRealBinLowN}{9}
\newcommand{\dmzRealBinLowMed}{0.59}
\newcommand{\dmzRealBinLowKsP}{0.80}
\newcommand{\dmzRealBinHighN}{27}
\newcommand{\dmzRealBinHighMedRange}{0.14--0.21}
\newcommand{\dmzRealTrendRho}{-0.31}
\newcommand{\dmzRealTrendP}{0.07}
\newcommand{\dmzRealNoChimeN}{15}
\newcommand{\dmzRealNoChimeKsP}{2\times10^{-5}}
\newcommand{\dmzRealChimeDsaN}{31}
\newcommand{\dmzRealChimeDsaKsP}{0.002}
\newcommand{\dmzRealChimeDsaCovNinetyFive}{0.87}
\newcommand{\dmzRealHaloLowKsP}{2\times10^{-6}}
\newcommand{\dmzRealHaloHighKsP}{5\times10^{-4}}
\newcommand{\dmzRealPruneShift}{6\times10^{-5}}
\newcommand{\dmzRealPlantOne}{0.948}
\newcommand{\dmzRealPlantThree}{0.958}
\newcommand{\dmzRealPlantTwo}{0.962}
\newcommand{\dmzRealFrbZ}{2.148}
\newcommand{\dmzRealFrbCalebLo}{1.63}
\newcommand{\dmzRealFrbCalebHi}{3.40}
\newcommand{\dmzRealSigmaN}{36}
\newcommand{\dmzRealSigmaCovSixtyEight}{0.56}
\newcommand{\dmzRealSigmaCovNinetyFive}{0.83}
\newcommand{\dmzRealSigmaMedPit}{0.19}
\newcommand{\dmzRealSigmaKsP}{2\times10^{-5}}
\newcommand{\dmzRealHostOneFiveN}{36}
\newcommand{\dmzRealHostOneFiveCovSixtyEight}{0.64}
\newcommand{\dmzRealHostOneFiveCovNinetyFive}{0.81}
\newcommand{\dmzRealHostOneFiveMedPit}{0.24}
\newcommand{\dmzRealHostOneFiveKsP}{0.004}
\newcommand{\dmzRealHostTwoN}{36}
\newcommand{\dmzRealHostTwoCovSixtyEight}{0.58}
\newcommand{\dmzRealHostTwoCovNinetyFive}{0.81}
\newcommand{\dmzRealHostTwoMedPit}{0.29}
\newcommand{\dmzRealHostTwoKsP}{0.02}
\newcommand{\dmzRealGoldN}{18}
\newcommand{\dmzRealGoldCovSixtyEight}{0.61}
\newcommand{\dmzRealGoldCovNinetyFive}{0.89}
\newcommand{\dmzRealGoldMedPit}{0.28}
\newcommand{\dmzRealGoldKsP}{0.14}
\newcommand{\dmzRealCovNinetyFiveK}{29}
\newcommand{\dmzRealCovNinetyFiveCiLo}{0.64}
\newcommand{\dmzRealCovNinetyFiveCiHi}{0.92}
\newcommand{\dmzRealCovNinetyFiveBinomP}{0.002}
\newcommand{\dmzRealSignP}{0.001}
\newcommand{\dmzRealNonChimeBelowHalf}{16}
\newcommand{\dmzRealNonChimeN}{16}
\newcommand{\dmzRealNonChimeSignP}{3\times10^{-5}}
\newcommand{\dmzRealNonChimeKsP}{1\times10^{-5}}
\newcommand{\dmzRealNonChimeCovNinetyFive}{0.69}
\newcommand{\dmzRealRatioMed}{1.85}
\newcommand{\dmzRealRatioIqrLo}{1.11}
\newcommand{\dmzRealRatioIqrHi}{3.32}
\newcommand{\dmzRealInSampleN}{39}
\newcommand{\dmzRealInSampleCovSixtyEight}{0.74}
\newcommand{\dmzRealInSampleCovNinetyFive}{0.95}
\newcommand{\dmzRealInSampleKsP}{0.02}
\newcommand{\dmzRealFirstCovSixtyEight}{0.53}
\newcommand{\dmzRealFirstCovNinetyFive}{0.78}
\newcommand{\dmzRealFirstPitTwoOneOneTwoFour}{0.015}
\newcommand{\dmzRealFirstPitTwoOneZeroFourOneZero}{0.0034}
\newcommand{\dmzRealNowPitTwoOneOneTwoFour}{0.165}
\newcommand{\dmzRealNowPitTwoOneZeroFourOneZero}{0.011}
\newcommand{\dmzSynN}{36}
\newcommand{\dmzSynM}{2000}
\newcommand{\dmzSynBandLoSixtyEight}{0.53}
\newcommand{\dmzSynBandHiSixtyEight}{0.83}
\newcommand{\dmzSynBandLoNinetyFive}{0.86}
\newcommand{\dmzSynFalseFail}{0.0445}
\newcommand{\dmzSynPowerHost}{0.12}
\newcommand{\dmzSynPowerF}{0.95}
\newcommand{\dmzSynPlantExpected}{0.9555}
\newcommand{\dmzSynDmzPowHost}{0.81}
\newcommand{\dmzSynDmzPowHostOneFive}{0.24}
\newcommand{\dmzSynDmzPowF}{0.014}
\newcommand{\dmzSynDmzPowMw}{0.08}

\newcommand{\code}[1]{\texttt{#1}}
\shorttitle{DM redshift intervals on localized FRBs}
\shortauthors{Barbere}
\begin{document}
\title{Dispersion-Measure Redshift Intervals Place Host-Identified Fast Radio Bursts Too Far Away}
\author[0009-0008-3289-4447]{Joseph Barbere}
\affiliation{Independent researcher}

\begin{abstract}
We tested the redshift posteriors $p(z\,|\,\mathrm{DM})$ that \code{zdm} infers from fast radio
burst (FRB) dispersion measures (DM), with the parameters of \citet{hoffmann2025}, on
\dmzRealNCert\ FRBs with spectroscopic hosts not used in that fit. The 95\% intervals contain
the host redshift for \dmzRealEOneCovNinetyFive\ of the bursts, below a floor of
\dmzSynBandLoNinetyFive\ fixed in advance. The failure is one-sided: \dmzRealEOneTailLo\ of the
probability integral transforms fall below the central 68\% band and \dmzRealEOneTailHi\ above
it ($p=\dmzRealEOneKsP$), and the posterior median is typically \dmzRealRatioMed\ times the true
redshift. It persists with the Galactic DM uncertainty of the fit included and with the mean host
DM raised by half. This sample cannot tell whether the population model, the selection of
identifiable hosts, or an unmodelled host DM is responsible.
\end{abstract}

\keywords{Radio transient sources (2008) --- Intergalactic medium (813) --- Astrostatistics
(1882) --- Redshift surveys (1378)}

\section{Test}
\label{sec:test}
For a burst at redshift $z$, the extragalactic DM is the sum of a cosmic term, whose mean and
scatter follow \citet{macquart2020}, and a host term redshifted by $1/(1+z)$. Inverting this
gives $p(z\,|\,\mathrm{DM})$. The \code{zdm} code \citep{james2022,zdm} adds each survey's
detection efficiency and an FRB population model. Its parameters were last fitted on bursts from
Parkes, ASKAP, DSA-110 and FAST \citep{hoffmann2025}. Point estimates of redshift from DM have
been compared with localized bursts before \citep{cordes2022}; here we ask whether the quoted
intervals contain the true redshift as often as they claim.

The sample is the localized-host compilation of the \code{FRB} repository \citep{frbrepo}:
\dmzRealNLocal\ bursts after duplicates are removed. Bursts that entered the
\citet{hoffmann2025} fit, \dmzRealNProd\ in all, were set aside: their 25 DSA-110 commissioning
bursts (their Section~2.3) and every CRAFT burst to the end of 2023 (their Section~2.4). We
required a spectroscopic redshift and a host association probability of at least 0.9. That
leaves \dmzRealNCert\ bursts: \dmzRealNChime\ from CHIME (\dmzRealGoldN\ from the gold sample of
\citealt{leung2025} and the repeaters FRB~20180916B and FRB~20201124A), \dmzRealNDsa\ from
DSA-110 \citep{sharma2024}, \dmzRealNAskap\ from ASKAP \citep{shannon2025} and \dmzRealNMeer\
from MeerKAT \citep{caleb2023}. Each burst was scored under the \code{zdm} survey model of its
discovery instrument. \citet{hoffmann2025} excluded CHIME, whose model needs repeaters to be
treated, and MeerKAT (their Section~2.1), so \dmzRealNChime\ of our bursts are scored with an
unfitted CHIME model summed over its declination bins, and one with an unfitted MeerKAT model.
The Milky Way DM is the NE2001 value \citep{cordes2002} plus 50~pc~cm$^{-3}$ for the halo. These
are the means of the fit, which also allowed uncertainties of half the NE2001 value and
15~pc~cm$^{-3}$ \citep[][Section~3.1]{hoffmann2025}; our primary run omits them.

We checked the discovery instrument and redshift of the non-CHIME bursts against the original
papers and corrected two after the first run. FRB~20201124A is listed as a MeerKAT burst but was
discovered by CHIME \citep{lanman2022}, with redshift 0.0979 \citep{fong2021}; FRB~20210410D has
DM 578.78~pc~cm$^{-3}$ and was found in the incoherent beam \citep{caleb2023}. Their PITs
(defined below) moved from \dmzRealFirstPitTwoOneOneTwoFour\ to \dmzRealNowPitTwoOneOneTwoFour\
and from \dmzRealFirstPitTwoOneZeroFourOneZero\ to \dmzRealNowPitTwoOneZeroFourOneZero, and the
\code{zdm} coverages from \dmzRealFirstCovSixtyEight\ and \dmzRealFirstCovNinetyFive\ to their
present values.

For each burst we computed the probability integral transform (PIT), the posterior probability
below the true redshift. A central $q$ interval contains the truth when the PIT lies between
$(1-q)/2$ and $(1+q)/2$. The rule was fixed before any real posterior was computed. The 68\%
and 95\% coverages must lie inside the central 95\% of their distributions for calibrated
samples of \dmzSynN\ bursts, which are \dmzSynBandLoSixtyEight--\dmzSynBandHiSixtyEight\ and at
least \dmzSynBandLoNinetyFive, and a Kolmogorov--Smirnov (KS) test of the PITs against a uniform
distribution must give $p>0.01$. On \dmzSynM\ calibrated synthetic samples the rule fails
\dmzSynFalseFail\ of the time. It has power \dmzSynPowerF\ against a doubled cosmic scatter but
only \dmzSynPowerHost\ against a doubled mean host DM. Because of that low power we changed the
rule after this control and before the real run: a failure is reported, and a pass would not
have been read as evidence of calibration. Planting each burst's redshift from its own posterior,
which makes the PITs uniform by construction, checks the code: the rule passes
\dmzRealPlantOne, \dmzRealPlantTwo\ and \dmzRealPlantThree\ of the time for the three
estimators, against \dmzSynPlantExpected\ expected.

We scored two other estimators on the same bursts. One is the \citet{macquart2020} posterior
with the same host and scatter parameters, a prior uniform in comoving volume, and no survey
selection. The other is \code{fruitbat} \citep{batten2019} with its default simulation-based
relation \citep{batten2021}, which includes no host term. Neither was fitted to these bursts,
and both were expected to fail.

\section{Results}
\label{sec:results}
All three estimators fail the rule (Figure~\ref{fig:pit}). For \code{zdm} the 68\% and 95\%
coverages are \dmzRealEOneCovSixtyEight\ and \dmzRealEOneCovNinetyFive. The 68\% coverage is in
band; the failure rests on the 95\% coverage, \dmzRealCovNinetyFiveK\ of \dmzRealNCert\ bursts
(binomial 95\% interval \dmzRealCovNinetyFiveCiLo--\dmzRealCovNinetyFiveCiHi;
$p=\dmzRealCovNinetyFiveBinomP$ if the true coverage were 95\%), and on the KS test. The median PIT
is \dmzRealEOneMedPit, and \dmzRealEOneNBelowHalf\ of the PITs lie below 0.5 (sign test
$p=\dmzRealSignP$). The other two estimators fail more strongly, with 95\% coverages of
\dmzRealETwoCovNinetyFive\ and \dmzRealEThreeCovNinetyFive. Since the PITs pile up low and almost
none lie in the upper tail, the intervals are displaced toward high redshift rather than too
narrow. The \code{zdm} posterior median is \dmzRealRatioMed\ times the true redshift (interquartile
range \dmzRealRatioIqrLo--\dmzRealRatioIqrHi). On the \dmzRealInSampleN\ bursts that entered the
fit, the coverages are \dmzRealInSampleCovSixtyEight\ and \dmzRealInSampleCovNinetyFive\
($p=\dmzRealInSampleKsP$).

The checks below were all chosen after the result was seen. All \dmzRealNonChimeN\ non-CHIME
bursts have PITs below 0.5 ($p=\dmzRealNonChimeSignP$); their 95\% coverage is
\dmzRealNonChimeCovNinetyFive\ and the KS test gives $p=\dmzRealNonChimeKsP$. Without the ASKAP and
MeerKAT bursts the 95\% coverage is \dmzRealChimeDsaCovNinetyFive, inside the band, but the KS
test still fails ($p=\dmzRealChimeDsaKsP$, \dmzRealChimeDsaN\ bursts). The \dmzRealGoldN\ gold
CHIME bursts alone show no failure (\dmzRealGoldCovSixtyEight\ and \dmzRealGoldCovNinetyFive,
$p=\dmzRealGoldKsP$), with little power at that size. With the Galactic DM uncertainty of the fit
the coverages are \dmzRealSigmaCovSixtyEight\ and \dmzRealSigmaCovNinetyFive\
($p=\dmzRealSigmaKsP$). A halo DM of 25 or 75~pc~cm$^{-3}$ gives $p=\dmzRealHaloLowKsP$ and
$\dmzRealHaloHighKsP$. All \dmzRealNDsa\ DSA-110 bursts and FRB~20180916B are listed in the
\code{zdm} survey files; removing them changes no PIT by more than $\dmzRealPruneShift$. A second
redshift for one burst and three DM values that differ between published tables do not change
the result.

The PIT of the DM at the known redshift does not depend on the redshift prior, though it does
depend on any selection acting on the host DM. It shows no failure: the coverages are
\dmzRealDmzOneCovSixtyEight\ and \dmzRealDmzOneCovNinetyFive\ and $p=\dmzRealDmzOneKsP$ for
\code{zdm}. In simulations at these redshifts this check detects a doubled mean host DM
\dmzSynDmzPowHost\ of the time, a mean host DM half again too large \dmzSynDmzPowHostOneFive,
a doubled cosmic scatter \dmzSynDmzPowF\ and an unmodelled 30~pc~cm$^{-3}$ \dmzSynDmzPowMw.
Raising the \code{zdm} mean host DM by factors of 1.5 and 2 leaves the 95\% coverage at
\dmzRealHostOneFiveCovNinetyFive\ and \dmzRealHostTwoCovNinetyFive\ ($p=\dmzRealHostOneFiveKsP$
and $\dmzRealHostTwoKsP$). A constant host-DM error up to a factor of 2 therefore does not
remove the failure. Of the explanations tested, the redshift distribution is the more likely
site: either the population model or the selection of identifiable hosts. For their DSA-110 commissioning sample, \citet{hoffmann2025}
used redshifts only below $\mathrm{DM}-\mathrm{DM_{ISM}}=183$~pc~cm$^{-3}$, because fainter hosts
at higher redshift are missed and a localized sample is then biased low in redshift at fixed DM.
Below that value our \dmzRealBinLowN\ bursts have a median PIT of \dmzRealBinLowMed\
($p=\dmzRealBinLowKsP$); above it, the \dmzRealBinHighN\ bursts have medians of
\dmzRealBinHighMedRange\ in three bins. That is the pattern such selection predicts, but
\dmzRealBinLowN\ bursts have little power and the trend of PIT with DM is not significant
($\rho=\dmzRealTrendRho$, $p=\dmzRealTrendP$). A host DM that rises with redshift would also bias
redshifts high \citep{masribas2026}; it was not tested.

FRB~20240304B \citep{caleb2026}, at $z=\dmzRealFrbZ$, was not in any statistic. The \code{zdm}
68\% interval, \dmzRealFrbOneLoSixtyEight--\dmzRealFrbOneHiSixtyEight, contains it; those of
the other two estimators (\dmzRealFrbTwoLoSixtyEight--\dmzRealFrbTwoHiSixtyEight\ and
\dmzRealFrbThreeLoSixtyEight--\dmzRealFrbThreeHiSixtyEight) do not. All three place it too far
away, as for the sample, though \citet{caleb2026} attribute part of its DM to a foreground
cluster.

\begin{figure}
\plotone{figures/dmzcal_pit.pdf}
\caption{Cumulative distributions of the PIT of the host redshift for the \dmzRealNCert\
bursts under the three estimators; a calibrated estimator follows the grey diagonal. Dotted:
the PIT of the DM at the known redshift under \code{zdm}.}
\label{fig:pit}
\end{figure}

\section{Limits}
\label{sec:limits}
The result applies to bursts with identified hosts, not to the unlocalized bursts for which
these estimators are intended. Either localized bursts differ in redshift at fixed DM from the
population the model describes, or the model's redshift distribution is wrong; this sample
cannot tell which. The three estimators were scored on the same bursts and are not independent
tests. NE2001 was not compared with other Galactic models, no explicit model of host
identification was fitted, and host magnitudes, available only for the \dmzRealGoldN\ gold CHIME bursts, were not used.

Reproducibility: \code{jansky\_research.dmzcal} and \code{scripts/dmzcal\_*.py}; every number
here is generated from \code{results/dmzcal\_*.json}. The plan, with the rule as fixed, and
the record of each step are \code{plans/99} and \code{survey/dmzcal-findings.md}.

\section*{AI use disclosure}
This paper was developed collaboratively by the author with Anthropic's Claude
\citep{anthropicclaude} (model Claude Opus~5.5, accessed 2026). Claude contributed to the
design, implementation, and testing of the \code{dmzcal} module and the drafting of this
manuscript. The author directed the work, independently reviewed and verified all code, results,
and citations, and takes full responsibility for the content. An AI/LLM is not an author.

\begin{acknowledgments}
We used the host compilation of the \code{FRB} repository and the published catalogues it
draws on; the analysis builds on the \code{jansky} project. The author acknowledges the
collaborative contribution of Anthropic's Claude.
\end{acknowledgments}

\software{\code{jansky-research} \citep{janskyresearch}, \code{zdm} \citep{zdm},
\code{fruitbat} \citep{batten2019}, \code{FRB} \citep{frbrepo}, Claude \citep{anthropicclaude},
Astropy, NumPy, SciPy, Matplotlib, jansky \citep{jansky}}

\bibliography{refs}
\end{document}
